\section{Threats to Validity and Limitations}
\label{sec:limitations}
\paragraph{Input and reference availability.}
Recovery is conditional on supplied code and candidate references. The completed experiments and planned primary evaluation use fixed reference inputs; they do not establish the accuracy or cost of automatic package identification and retrieval. Missing dependencies, unavailable package variants, patched library copies, and uncaptured chunks can limit the available correspondences. The planned verification dataset records missing or incompatible references as preparation exclusions rather than silently repairing inputs or attributing these failures to the matcher. The RQ4 practical-use protocol remains undefined, and no downstream improvement is established in this draft.

\paragraph{Deliberate obfuscation.}
Our evaluation targets bundling and minification and does not establish robustness to deliberate JavaScript obfuscation. Prior web measurements found explicit obfuscation substantially less prevalent than minification~\cite{skolka2019anything}; these historical findings do not establish its current prevalence. Property-name concealment or renaming can invalidate POG anchors, while control-flow flattening and inserted dispatchers or wrappers can disrupt relationship graphs and candidate comparison. Packing, runtime-generated code, and virtualization may hide or replace the original function bodies, limiting the available static correspondences. We leave systematic evaluation under these transformations and extensions for recovering obfuscated functions to future work; the reported recovery and reachability results should not be generalized to such inputs.

\paragraph{Source identity versus semantic equivalence.}
Our objective is correspondence to a particular source definition. Different definitions can implement the same behavior, so a valid semantic acceptance can still disagree with this identity oracle. Conversely, a source-linked bundled function may be specialized to its actual callers and differ on unreachable inputs. We therefore describe precision as source-identity agreement, not as the probability that a semantic-equivalence proof is correct. We do not treat every identity error as an independently established semantic false positive.

\paragraph{External bindings and missing program context.}
Function bodies alone may not identify decisive enclosing bindings. Eight of the 14 wrong identities in the 27B run involve different regex/replacement helpers or public/private data stores behind nearly identical wrapper bodies. Minifying the source improves syntactic comparability but does not recover the bundled binding's origin. The permissive prompt allows inferred correspondences and accepts without a demonstrated difference. This policy can accept a plausible but wrong binding mapping. The limitation is not the mere presence of an external variable: sufficient local differences can still justify rejection. Additional binding evidence or relevant helper definitions would be needed to resolve ambiguous cases reliably.

\paragraph{Permissive validation and function boundaries.}
The adopted prompt explicitly excludes corner cases and side effects, weakening the comparison relative to full JavaScript semantics. This restriction can hide differences in data writes or helper calls, although T/F outputs do not reveal whether it caused any particular acceptance. Two wrong setter matches have different arguments and enclosing roles; another confuses a layout-test result with the jQuery object. Three more confuse enclosing functions with nested functions. Shared nested code must not establish equivalence of the root functions. Neither a larger model nor a minified reference guarantees correct boundary or behavior analysis.

\paragraph{Validator-dependent acceptance policy.}
The same permissive prompt yields accepted-pair assessed precision of 94.65\% with 27B and 61.03\% with 4B. The absence of observed correct-pair rejections in the 4B run does not make it a reliable validator: false acceptances can terminate a search early or block other proposals. First-positive stopping and one-to-one allocation can amplify this difference. We observe the decisions and their pipeline consequences, but T/F-only outputs do not establish whether model capacity, the default-to-accept instruction, omitted side-effect checks, or another factor caused a particular error. The comparison does not establish that smaller validators are generally unsuitable; another prompt or validation policy remains untested.

\paragraph{Propagation and irreversible allocation.}
An incorrect match can change later neighborhoods and consume another function's source identity. Five experiment~4 cases and 13 experiment~5 cases have earlier wrong owners that exclude the correct source from a later candidate pool. These are observed allocation consequences, not counterfactual gains from rollback. First-positive checking does not revisit accepted allocations or inspect later candidates. Conflicts block an accepted correct pair for 17 supported targets with 27B and 92 with 4B. A positive judgment is therefore not necessarily a recovered identity.

\paragraph{Source-map labels and unresolved functions.}
Labels depend on mapping density and the voting rule. Tiny, inlined, merged, or otherwise transformed functions can lack a supported identity. We retain 45 unassessed promotions in experiment~4 and 38 in experiment~5 separately from correct and wrong predictions. Independent own-body mapping witnesses agree with all 14 experiment~4 and 62 experiment~5 wrong-identity labels, strengthening those specific labels without validating every source-map segment or proving equivalence. A missing identity is not evidence that the bundler created a new library function.

\paragraph{Static structure and isolated minification.}
The graph is an incomplete static approximation, supplemented by containment and source order. Proximity supplies candidates rather than proving semantic dependence. A missing intermediate match can block one-hop expansion, while a shared factory can expose many unrelated candidates. Isolated-function minification lacks enclosing constants, bindings, and compiler context, so it need not reproduce the deployed bundler's transformations. Input-size restrictions and one-to-one matching also limit support for large functions, duplicates, inlining, merging, and specialization.

\paragraph{Experimental scope and comparison.}
The pilot uses one known library version and one crafted application that exposes the whole library, with minification enabled and tree shaking disabled. Prompt development and diagnostics used this subject; there is no held-out library or version for the model-driven comparison. The new-repository graph experiment is held out from analyzer development, but does not run the models. Repeated function occurrences across builds, repeated queries, and cached decisions are correlated. Experiment~4 changes source preprocessing, the embedding instruction, validator input and prompt, and the validation budget together. Its improvement over experiment~3 cannot be attributed to one component. The experiment~5 comparison holds the pipeline, prompt, inputs, and embedding vectors fixed, but compares deployed models rather than parameter count in isolation. Both use NVFP4 weights; quantization may affect the models differently. Their shared-pair subset is selected by both trajectories, and the 4B run does not have complete individual semantic adjudication of every wrong identity. All timings are single observed runs with fresh experiment caches and within-run reuse, not repeated, hardware-normalized benchmarks. These model comparisons do not establish complete recovery, a zero-false-positive guarantee, or model generalization to other libraries and optimization settings.

\paragraph{Connectivity, candidate breadth, and bundler coverage.}
The graph is designed to admit useful candidates, including approximate references and proximity relations. Extra edges do not constitute accepted identities: the model stage makes that decision. The oracle simulation selects only supported true candidates and bypasses ranking and validation, so it cannot quantify the cost of larger candidate sets, top-five competition, or model mistakes on those sets. Its rates are conditional on the available source-map labels, parseable source inventory, confirmed seed policy, and promotion restrictions. Zero-supported cases and same-chunk ambiguities remain visible. The getter/reference audit characterizes approximate candidate edges; we retain the unchanged graph rather than treating edge precision as the optimization objective of this stage.

The survey motivates the selected bundler families but is self-selected developer usage data, not a deployment census. Application selection depends on available pinned builds, and the five new repositories share some library families with development subjects. The large development gain is concentrated in one Rolldown icon library; the Parcel holdout has only 14 targets. The Turbopack holdout uses an explicit backend change and a font-CSS fixture. Runtime-only context completeness, plugin-generated wrappers, development/HMR protocols, module federation, dynamically computed assets, and future protocol changes remain possible coverage boundaries. These limits do not negate the measured gains, but prevent an ecosystem-wide success claim or a claim that every bundler version is supported.
